\documentclass[spanish,letterpaper,oneside]{article}

\def\documenttitle {Foro de semana 1: Derecho penal especial mexicano}
\def\documentsubtitle {Participacion y respuestas del foro diagnostico}
\def\documentsubject {Semana 1 - Derecho penal especial mexicano}
\def\documentauthor {Martin Jonathan de la Cruz}
\def\coursename {Derecho penal especial mexicano}
\def\coursecode {020}
\def\universityname {Universidad Abierta y a Distancia de Mexico}
\def\universityfaculty {Licenciatura en Derecho}
\def\universitydepartment {Derecho penal especial mexicano}
\def\universitydepartmentimage {departamentos/UnADM}
\def\universitydepartmentimagecfg {height=1.57cm}
\def\universitylocation {Roma Norte, Ciudad de Mexico}
\def\authortable {
  \begin{tabular}{ll}
    \textbf{Alumno:} & Martin Jonathan de la Cruz \\
    Matricula: & ES2611202040 \\
    Figura docente: & Veronica Ortiz Lopez \\
    Semestre/Bloque: & 2 / 2 \\
    Estado: & Borrador para revision \\
    \multicolumn{2}{l}{Fecha de realizacion: \today} \\
    \multicolumn{2}{l}{\universitylocation}
  \end{tabular}
}

\input{template}
\usepackage{helvet}
\renewcommand{\familydefault}{\sfdefault}
\usepackage{array}
\usepackage{booktabs}
\usepackage{longtable}
\usepackage{calc}
\usepackage{tikz}
\usetikzlibrary{arrows.meta,positioning,calc}
\setcitestyle{authoryear,open={(},close={)}}
\providecommand{\tightlist}{\setlength{\itemsep}{0pt}\setlength{\parskip}{0pt}}
\providecommand{\pandocbounded}[1]{#1}
\providecommand{\passthrough}[1]{#1}
\setlength{\emergencystretch}{3em}

\usepackage[most]{tcolorbox}
\definecolor{unadmGreenDark}{HTML}{174A3A}
\definecolor{unadmGreen}{HTML}{5F8F3A}
\definecolor{unadmPaper}{HTML}{F6F7F2}
\newtcolorbox{forobox}[1][]{%
  enhanced, breakable,
  colback=unadmPaper, colframe=unadmGreen,
  boxrule=0pt, leftrule=3pt, arc=1.2mm,
  left=4mm, right=4mm, top=3mm, bottom=3mm,
  fonttitle=\bfseries\color{white}, coltitle=white, colbacktitle=unadmGreenDark,
  attach boxed title to top left={xshift=4mm,yshift=-2mm},
  boxed title style={colframe=unadmGreenDark,arc=0.6mm},
  #1
}

\def\coverwatermarkenabled {true}
\def\coverwatermarkimage {img/departamentos/UnADM.pdf}
\def\coverwatermarkopacity {0.16}
\def\coverwatermarkwidth {0.72\paperwidth}
\def\coverwatermarkxshift {0cm}
\def\coverwatermarkyshift {-0.35cm}
\newcommand{\insertcoverwatermark}{%
  \ifthenelse{\equal{\coverwatermarkenabled}{true}}{%
    \AddToShipoutPictureBG*{%
      \begin{tikzpicture}[remember picture,overlay]
        \node[opacity=\coverwatermarkopacity,inner sep=0pt,
          xshift=\coverwatermarkxshift,yshift=\coverwatermarkyshift]
          at (current page.center) {%
            \includegraphics[width=\coverwatermarkwidth]{\coverwatermarkimage}%
          };
      \end{tikzpicture}%
    }%
  }{}%
}

\begin{document}
\insertcoverwatermark
\templatePortrait
\templatePagecfg
\onehalfspacing
\templateIndex
\templateFinalcfg

\section{Identificacion y estado}

Propuesta de participacion basada en el proyecto, elaborada con apoyo de GitHub Copilot. Requiere revision personal; no publicada. No acredita conocimientos previos ni experiencias personales confirmadas. Edad y residencia, cuando se solicitan, permanecen pendientes.

\section{Participacion principal}
\begin{forobox}[title={Participacion principal -- propuesta no publicada}]

\textbf{1. Cuando escuchas hablar de Derecho Penal, ¿que es lo primero
que relacionas con esta area del Derecho?}

Lo relaciono con las conductas consideradas delitos, la proteccion de
las personas y los limites que debe respetar el Estado al investigar y
sancionar. Tambien pienso en la importancia de distinguir una acusacion
de una responsabilidad demostrada, porque una respuesta penal afecta
derechos y no debe basarse solamente en sospechas.

\textbf{2. ¿Has tenido algun acercamiento previo con temas relacionados
con el Derecho Penal?}

Mi acercamiento dentro de este proyecto ha sido academico y documental:
he trabajado temas de garantias constitucionales, derechos humanos y
analisis juridico. Desde ese enfoque me interesa comprender como se
protegen los derechos frente a las decisiones de la autoridad. Ese
trabajo no equivale a experiencia profesional en asuntos penales.

\textbf{3. Desde tu punto de vista, ¿para que consideras que puede
servirte el estudio del Derecho Penal en tu formacion y futuro ejercicio
profesional?}

Puede ayudarme a distinguir los hechos relevantes de las suposiciones,
comprender que debe analizarse antes de afirmar que existe un delito y
reconocer los derechos de las personas involucradas. Espero fortalecer
mi argumentacion y aprender a orientar con responsabilidad, sin prometer
resultados ni emitir conclusiones apresuradas.

\textbf{4. ¿Te gustaria desarrollarte profesionalmente en algun area
relacionada con el Derecho Penal? ¿Por que?}

Me interesa explorar su relacion con la proteccion de derechos humanos y
las garantias de las personas acusadas y de las victimas. No lo planteo
como una especializacion ya decidida, sino como un campo que deseo
conocer mejor. Considero importante que cualquier intervencion juridica
combine analisis de los hechos, respeto a los derechos y responsabilidad
profesional.

\textbf{5. Una persona te comenta que ha sido acusada de cometer un
delito y te pide orientacion. Sin necesidad de utilizar terminos
juridicos, ¿que informacion considerarias importante conocer antes de
orientarla?}

Preguntaria que sucedio, cuando y donde, que acusacion le comunicaron y
si ha recibido una citacion o esta detenida. Tambien buscaria saber si
cuenta con una persona abogada que la asista y si existen documentos o
testigos. Evitaria asumir que es culpable o prometer un resultado;
primero es necesario comprender su situacion y asegurar que reciba
orientacion adecuada.

\textbf{6. Si una persona te dijera que ha sido victima de un delito,
¿que informacion buscarias conocer y que considerarias importante hacer
para poder orientarla?}

Primero preguntaria si esta a salvo y si necesita atencion medica o
apoyo inmediato. Despues conoceria lo ocurrido, el lugar, la fecha y si
hay documentos o testigos. Procuraria orientarla hacia ayuda profesional
y las autoridades correspondientes, respetando su privacidad y evitando
que tenga que enfrentar por su cuenta a quien le causo el dano.

\textbf{7. Reflexion final}

Considero que el estudio del Derecho Penal Especial Mexicano puede
contribuir a mi formacion profesional porque me permitira analizar con
mayor cuidado las conductas y sus consecuencias, distinguir el objeto de
una acusacion y fundamentar mejor mis argumentos. Quiero relacionar ese
aprendizaje con las garantias constitucionales y la proteccion de
derechos, para que la orientacion juridica no se reduzca a buscar una
sancion, sino que atienda la situacion de todas las personas
involucradas.

\end{forobox}

\section{Respuestas a participaciones}

\begin{forobox}[title={Respuesta 1 -- pendiente de intervencion real}]

Pendiente: seleccionar y leer una participacion real, identificar su idea central y redactar una retroalimentacion respetuosa y especifica. No se simula una respuesta publicada ni se atribuye una experiencia de otra persona al estudiante.

\end{forobox}

\begin{forobox}[title={Respuesta 2 -- pendiente de intervencion real}]

Pendiente: seleccionar y leer una participacion real, identificar su idea central y redactar una retroalimentacion respetuosa y especifica. No se simula una respuesta publicada ni se atribuye una experiencia de otra persona al estudiante.

\end{forobox}

\end{document}